\subsection{内射模}

定义2. --- 称A-模E为内射的，如果对每个由A-模及同态构成的正合序列

\[
\mathbf {M} ^ {\prime} \xrightarrow {u} \mathbf {M} \xrightarrow {v} \mathbf {M} ^ {\prime \prime},\tag{19}
\]

由Z-线性映射构成的序列

\[
\operatorname{Hom} _ {\mathbf {A}} \left(\mathbf {M} ^ {\prime \prime}, \mathbf {E}\right) \xrightarrow {\operatorname{Hom} _ {\mathbf {A}} (v , 1)} \operatorname{Hom} _ {\mathbf {A}} (\mathbf {M}, \mathbf {E}) \xrightarrow {\operatorname{Hom} _ {\mathbf {A}} (u , 1)} \operatorname{Hom} _ {\mathbf {A}} \left(\mathbf {M} ^ {\prime}, \mathbf {E}\right)\tag{20}
\]

是正合的。

引理3. --- 对于A-模E为内射的，当且仅当对每个单射的 A-线性映射 u : M' → M，映射

\[
\operatorname{Hom} _ {\mathbf {A}} (u, 1): \operatorname{Hom} _ {\mathbf {A}} (\mathrm{M}, \mathrm{E}) \rightarrow \operatorname{Hom} _ {\mathbf {A}} (\mathrm{M} ^ {\prime}, \mathrm{E})
\]

是满射。

若E是内射的且若 \(u: \mathbf{M}' \to \mathbf{M}\) 是单射的，则序列 \(0 \to \mathbf{M}' \xrightarrow{u} \mathbf{M}\) 是正合的，因此序列 Hom (M, E) \(\xrightarrow{\text{Hom}(u,1)}\) Hom ( \(\mathbf{M}'\) , E) \(\to 0\) 是正合的，且 Hom ( \(u,1\) ) 是满射。反之，考虑正合序列 (19)；设 \(\mathbf{M}_{1}'' = v(\mathbf{M})\) 并设 \(i: \mathbf{M}_{1}'' \to \mathbf{M}''\) 为典范单射， \(p: \mathbf{M} \to \mathbf{M}_{1}''\) 为映射 \(m \mapsto v(m)\) 。序列 \(\mathbf{M}' \xrightarrow{u} \mathbf{M} \xrightarrow{p} \mathbf{M}_{1}'' \to 0\) 是正合的；由 II，第～36 页，定理 1，序列

\[
\operatorname{Hom} _ {\mathbf {A}} \left(\mathrm{M} _ {1} ^ {\prime \prime}, \mathrm{E}\right) \xrightarrow {\operatorname{Hom} (p , 1)} \operatorname{Hom} _ {\mathbf {A}} (\mathrm{M}, \mathrm{E}) \xrightarrow {\operatorname{Hom} (u , 1)} \operatorname{Hom} _ {\mathbf {A}} \left(\mathrm{M} ^ {\prime}, \mathrm{E}\right)
\]

是正合的。此外，我们有 Hom \((v, 1) = \text{Hom } (p, 1) \circ \text{Hom } (i, 1)\) 。若 E 满足引理的条件，则 Hom \((i, 1)\) 是满射，因此 Hom \((v, 1)\) 的像也是 Hom \((p, 1)\) 的像，且序列 (20) 是正合的。

注记. --- 设 E 是一个内射 A-模， \(u: \mathbf{M}' \to \mathbf{M}\) 和 \(f: \mathbf{M}' \to \mathbf{E}\) 为 A-模的同态。若 \(\operatorname{Ker} u \subset \operatorname{Ker} f\) ，则存在一个同态 \(g: \mathbf{M} \to \mathbf{E}\) 使得 \(g \circ u = f\) 。这实际上由前述内容推出，将其应用于单射同态 \(\mathbf{M}' / \operatorname{Ker} u \to \mathbf{M}\) （由 \(u\) 得出）。

命题 9. --- 设 \((\mathbf{E}_i)_{i\in I}\) 是 A-模的一个族， \(\mathbf{E} = \prod \mathbf{E}_i\) 它们的积。对于 A-模 E 为内射的，必要且充分条件是每个 \(\mathbf{E}_i\) 都是内射的。

设 \(u: \mathbf{M}' \to \mathbf{M}\) 是 A-模的一个内射同态。对于积同态 \(\prod_{i \in I} \operatorname{Hom}_{\mathbf{A}}(\mathbf{M}, \mathbf{E}_i) \to \prod_{i \in I} \operatorname{Hom}_{\mathbf{A}}(\mathbf{M}', \mathbf{E}_i)\) 为满射的，必要且充分条件是每个同态 \(\operatorname{Hom}_{\mathbf{A}}(\mathbf{M}, \mathbf{E}_i) \to \operatorname{Hom}_{\mathbf{A}}(\mathbf{M}', \mathbf{E}_i)\) 都是满射的（II，第～10 页，命题 5）；这证明了该命题，因为 \(\prod_{i \in I} \operatorname{Hom}_{\mathbf{A}}(\mathbf{M}, \mathbf{E}_i)\) 典范地等同于 \(\operatorname{Hom}_{\mathbf{A}}(\mathbf{M}, \mathbf{E})\) 。

命题 10. --- 设 E 是一个 A-模。对于 E 为内射的，必要且充分条件是，对于 A 的每个理想 a 和每个 A-同态 \(f: \mathfrak{a} \to \mathbb{E}\) ，存在 \(e \in \mathbb{E}\) ，使得 \(f(a) = ae\) 对所有 \(a \in \mathfrak{a}\) 成立。

假设 E 是内射的；设 a 是 A 的一个理想， \(f : a \to E\) 为一个 A-同态，并设 \(i : a \to A\) 为典范单射。则映射

\[
\operatorname{Hom} _ {\mathbf {A}} (i, 1): \operatorname{Hom} _ {\mathbf {A}} (\mathrm{A}, \mathrm{E}) \rightarrow \operatorname{Hom} _ {\mathbf {A}} (\mathfrak {a}, \mathrm{E})
\]

是满射（定义 2）；若 \(g \in \operatorname{Hom}_{\mathbf{A}}(\mathbf{A}, \mathbf{E})\) 满足 \(f = g \circ i\) ，我们有

\[
f (a) = g (a) = a g (1)
\]

对所有 \(a \in \mathfrak{a}\) 成立。

反之，假设命题中的条件成立；设 M 为一个 A-模，N 为 M 的一个子模， \(u: \mathbf{N} \to \mathbf{E}\) 一个A-同态，让我们证明存在一个A-同态 \(\overline{u}: \mathbf{M} \to \mathbf{E}\) 扩展 \(u\) （参见引理3）。设 \(\mathcal{P}\) 为二元组的集合， \((\mathbf{P}, v)\) 其中P是M的包含N的子模，且 \(v\) 是从 P 到 E 的同态并扩张 \(u\) ；集合 \(\mathcal{P}\) 按扩张关系排序是归纳的：如果 \((\mathbf{P}_{j}, v_{j})\) 是 P 的元素的全序族，令 \(Q = \cup P_{j}\) 并设 \(w : Q \to E\) 是唯一诱导出各映射 \(v_{j}\) （在对应的 \(P_{j}\) 上，对所有 j 成立）的映射；则 \((\mathbf{Q}, w) \in \mathcal{P}\) 且 \((\mathbf{Q}, w)\) 优超 \((\mathbf{P}_{j}, v_{j})\) 对所有 j 成立。然后令 \((\mathbf{P}, v)\) 为P的一个极大元（E，III，第20页，定理2）；只需证明P = M。设 \(x \in M\) ，并设 a 为由那些 \(a \in A\) （满足 \(ax \in P\) ）组成的理想；令 \(f(a) = v(ax)\) ，其中 \(a \in a\) ；于是得到一个 A-同态 \(f : a \to E\) 。设 e 为 E 中使得 \(f(a) = ae\) 对所有 \(a \in a\) 都成立的元素。令 \(P' = P + Ax\) 并设 \(v' : P' \to E\) 是唯一的A-同态，使得 \(v'(p + ax) = v(p) + ae\) （其中 \(p \in P, a \in A\) ）；然后 \((\mathbf{P}', v')\) 属于 P 且优超 \((\mathbf{P}, v)\) ，因此 \(P' = P\) ，即 \(x \in P\) ，这便完成了证明。

推论1. --- 若环A是左Noether环，则任何由内射A-模的直和构成的模都是内射的。

设 \((\mathrm{E}_{i})_{i\in\mathbb{I}}\) 是内射 A-模的一个族，设 E 为它们的直和，a 为 A 的一个理想，且 \(u:a\to E\) 为一个 A-同态。由于 A 是诺特环，a 是有限生成的，因此典范映射

\[
\varphi : \bigoplus_ {i \in I} \operatorname{Hom} _ {\mathbf {A}} (\mathfrak {a}, E _ {i}) \rightarrow \operatorname{Hom} _ {\mathbf {A}} (\mathfrak {a}, E)
\]

是双射；设 \((u_{i})\) 为 u 在 \(\varphi\) 下的原像。由于每个 \(E_{i}\) 是内射的，且族 \((u_{i})\) 具有有限支撑，存在E中的一个元素 \((e_{i})_{i\in\mathbf{I}}\) 使得 \(u_{i}(a)=ae_{i}\) 对所有 \(a\in\mathfrak{a}\) 且对每个 \(i\in I\) ，因此 \(u(a)=a((e_{i}))\) 对所有 \(a\in\mathfrak{a}\) ，且E是内射的。

注记.—— 若每个作为内射 A-模直和的 A-模都是内射的，则环 A 是左 Noether 的（X，第 170 页，习题 21）。

设 A 是一个整环。称 A-模 E 是可除的，如果位似 \(a_{E}\) 对环 A 的每个非零元素 \(a\) 都是满射的。

推论2. —— 假设A是一个整环。

\begin{enumerate}
\def\labelenumi{\alph{enumi})}
\item
  每个内射A-模都是可除的。
\item
  每一个无挠（II，第115页）且可除的A-模都是内射的。
\item
  若 A 是主理想整环，则每个可除 A-模都是内射模。
\end{enumerate}

若 \(a \in A\) 非零，则 \(a_{A}\) 是单射；另一方面，对于每个A-模E，位似 \(a_{E}\) 典范地等同于

\[
\operatorname{Hom} \left(a _ {\mathrm{A}}, 1\right): \operatorname{Hom} _ {\mathrm{A}} (\mathrm{A}, \mathrm{E}) \rightarrow \operatorname{Hom} _ {\mathrm{A}} (\mathrm{A}, \mathrm{E}),
\]

因此，E 是可除的当且仅当 Hom \((a_{\mathbf{E}}, 1_{\mathbf{E}})\) 对所有非零的 \(a \in \mathbf{A}\) 都是满射。断言 \(a)\) 由此由定义 2（X，第 15 页）得出。

设 E 是一个可除 A-模；假设 A 是主理想整环（相应地，E 是无挠的），并且让我们通过应用命题 10 来证明 E 是内射的。设 a 是 A 的一个理想，且 \(f: \mathfrak{a} \to \mathbf{E}\) 是一个 A-同态。设 \(x \in \mathfrak{a}\) 使得 \(\mathfrak{a} = \mathbf{A}x\) （相应地，使得 \(x \neq 0\) 如果 \(\mathfrak{a} \neq 0\) ），并设 \(e \in \mathbf{E}\) 使得 \(xe = f(x)\) 。让我们证明对每一个 \(a \in \mathfrak{a}\) ，我们有

\[
f (a) = a e;
\]

这是显然的，如果 \(a \in Ax\) ，因此当 A 是主理想整环时断言成立；如果 E 是无挠的，则对 \(xa \in \mathfrak{a}\) ，我们有 \(xf(a) = f(ax) = axe\) ，因此 \(f(a) = ae\) ，因为 \(x\) 是非零的，如果 \(a \neq 0\) 。

如果A是一个整环，则A的分式域K是一个内射A-模。如果A是主理想整环，则K/A是一个内射A-模。

\begin{enumerate}
\def\labelenumi{\arabic{enumi})}
\setcounter{enumi}{1}
\item
  例如，Z-模Q和Q/Z是内射的。
\item
  设 A 为主理想环，且设 \(\iota\) 为 A 的非零元素。则 A/aA 是 A/aA-模的内射模（X，第 170 页，习题 20）。
\end{enumerate}

